\begin{table*}[tp]
\centering
\renewcommand{\arraystretch}{1.3}
\caption{Requisitos No Funcionales del Sistema -- Grupo G5}
\label{tab:rnf_proyectofinal}
\begin{tabularx}{\textwidth}{|c|X|c|c|}
\hline
\textbf{ID} & \textbf{Descripción} & \textbf{Categoría} & \textbf{Fuente} \\
\hline
RNF-01 & Firmware en VS Code con PlatformIO. Prohibido .ino y Arduino IDE. & Restricción & Enunciado X \\
\hline
RNF-02 & Código en GitHub con historial de commits organizado que evidencie avance progresivo. & Trazabilidad & Enunciado X \\
\hline
RNF-03 & Operación continua mínimo 30 min sin fallos de sensor, actuador ni desconexión IoT. & Disponibilidad & Enunciado XXI \\
\hline
RNF-04 & Prototipo portable, peso estimado menor a 5 kg en operación completa. & Portabilidad & Acta 1 núm. 3 \\
\hline
RNF-05 & Latencia máxima de 15 s desde adquisición hasta actualización de canal en ThingSpeak (restricción free-tier). & Rendimiento & ThingSpeak API \\
\hline
RNF-06 & Contenedor sin fugas durante pruebas de riego; con sistema de drenaje. & Fiabilidad & Acta 1 núm. 6 \\
\hline
RNF-07 & ESP32 y conexiones protegidas de salpicaduras durante operación del riego. & Seg. eléctrica & Diseño propio \\
\hline
RNF-08 & Alimentación 5 V DC USB o power bank para demostración en laboratorio. & Energía & Diseño propio \\
\hline
RNF-09 & El techo móvil debe abrir y cerrar en menos de 3 s desde la activación del servo. & Rendimiento & Diseño propio \\
\hline
RNF-10 & Documentación de evidencias en bitácora técnica con fecha y responsables. & Documentación & Enunciado XVII-D \\
\hline
\end{tabularx}
\end{table*}